\DocumentMetadata{
  pdfversion=2.0,
  pdfstandard=ua-2,
  lang=en-US,
  tagging=on,
  tagging-setup={math/setup=mathml-SE,tabsorder=structure}
}
\documentclass[11pt,letterpaper]{article}
\usepackage[margin=0.68in,includefoot]{geometry}
\usepackage{newcomputermodern}
\usepackage{amsmath,amscd,mathtools}
\usepackage{tikz,adjustbox}
\providecommand{\square}{\mdlgwhtsquare}
\usepackage[hidelinks]{hyperref}
\hypersetup{
  pdftitle={Algebra First-Year Practice Exam B, Part 2},
  pdfauthor={University of Florida Department of Mathematics},
  pdfsubject={Graduate mathematics examination},
  pdfdisplaydoctitle=true,
  bookmarksopen=true
}
\setcounter{secnumdepth}{0}
\setlength{\parindent}{0pt}
\setlength{\parskip}{0.28em}
\setlength{\emergencystretch}{3em}
\setlength{\leftmargini}{2.15em}
\setlength{\leftmarginii}{2.65em}
\setlength{\leftmarginiii}{3em}
\pagestyle{empty}
\raggedbottom
\allowdisplaybreaks
\NewTaggingSocketPlug{block/list/label}{examproblem}{%
  \tagstructbegin{tag=Lbl}\tagstructend%
  \tagstructbegin{tag=\LBody}%
  \tagstructbegin{tag=\ProblemHeadingTag,title-o={\ProblemHeadingText}}%
  \tagmcbegin{tag=\ProblemHeadingTag,actualtext={\ProblemHeadingText}}%
  #2%
  \tagmcend%
  \tagstructend%
}
\newenvironment{examproblems}[2]{%
  \def\ProblemHeadingTag{#2}%
  \AssignTaggingSocketPlug{block/list/label}{examproblem}%
  \begin{enumerate}%
  \setcounter{enumi}{#1}%
  \renewcommand{\labelenumi}{\textbf{\theenumi.}}%
  \setlength{\topsep}{0.35em}%
  \setlength{\partopsep}{0pt}%
  \setlength{\itemsep}{0.48em}%
  \setlength{\parsep}{0.18em}%
}{\end{enumerate}}
\newenvironment{examsubparts}{%
  \AssignTaggingSocketPlug{block/list/label}{default}%
  \begin{enumerate}%
  \setlength{\topsep}{0.18em}%
  \setlength{\partopsep}{0pt}%
  \setlength{\itemsep}{0.16em}%
  \setlength{\parsep}{0.1em}%
}{\end{enumerate}}
\newenvironment{examinstructions}{%
  \par\begingroup\itshape
}{%
  \par\endgroup\vspace{0.18em}
}
\makeatletter
\renewcommand\subsection{\@startsection{subsection}{2}{0pt}%
  {-1.25ex plus -.25ex minus -.1ex}%
  {0.55ex plus .1ex}%
  {\normalfont\large\bfseries}}
\renewcommand{\@maketitle}{%
  \newpage\null\vskip -1em%
  {\footnotesize\bfseries
    \href{https://www.ufl.edu/}{University of Florida}, \href{https://math.ufl.edu/}{Department of Mathematics}\par}%
  \vskip -0.5em%
  \tagstructbegin{tag=H1,title={Algebra First-Year Practice Exam B, Part 2}}%
  \tagpdfparaOff%
  \tagmcbegin{tag=H1}%
    {\LARGE\bfseries\href{https://gma.math.ufl.edu/exams/algebra-first-year/}{Algebra First-Year Practice Exam B}, Part 2\par}%
  \tagmcend%
  \tagpdfparaOn%
  \tagstructend%
  \vskip 0.25em%
  \tagmcbegin{artifact}%
    \hrule height 1pt%
  \tagmcend%
  \vskip 1em%
}
\makeatother
\title{Algebra First-Year Practice Exam B, Part 2}
\author{}
\date{}
\begin{document}
\maketitle
\thispagestyle{empty}
\begin{examinstructions}
Answer four problems. Write your answers clearly in complete English sentences. You may quote results (within reason) as long as you state them clearly.
\end{examinstructions}
\begin{examproblems}{0}{H2}
\def\ProblemHeadingText{Problem 1.}
\item
This problem has three parts.
\begin{examsubparts}
\item[\textnormal{(a)}]
Prove that a prime \(p \in \mathbb{N}\) such that \(p \equiv 3 \pmod 4\) is prime when considered as an element in the ring of Gaussian integers \(\mathbb{Z}[i]\).
\item[\textnormal{(b)}]
Prove that a prime \(p \in \mathbb{N}\) such that \(p \equiv 1 \pmod 4\) is a product of two non-associate primes in \(\mathbb{Z}[i]\).
\item[\textnormal{(c)}]
Give a factorization of \(30\) as the product of a unit and powers of non-associate primes in \(\mathbb{Z}[i]\).
\end{examsubparts}
\def\ProblemHeadingText{Problem 2.}
\item
Let~\(R\) be the ring of \(n \times n\) matrices over a field~\(F\), where \(n \geq 2\).
\begin{examsubparts}
\item[\textnormal{(a)}]
Show that~\(R\) has no ideals other than~\(0\) and~\(R\).
\item[\textnormal{(b)}]
Exhibit a nonzero proper left ideal of~\(R\).
\end{examsubparts}
\def\ProblemHeadingText{Problem 3.}
\item
Find one representative of each conjugacy class of elements of order~\(4\) in the group \(GL(6,\mathbb{Q})\).
\def\ProblemHeadingText{Problem 4.}
\item
Let~\(F\) be a field.
\begin{examsubparts}
\item[\textnormal{(a)}]
What does it mean to say that a field extension of~\(F\) is algebraic?
\item[\textnormal{(b)}]
Let \(F \subset K \subset L\) be three fields. Prove that if~\(K\) is algebraic over~\(F\) and~\(L\) is algebraic over~\(K\), then~\(L\) is algebraic over~\(F\).
\end{examsubparts}
\def\ProblemHeadingText{Problem 5.}
\item
A commutative ring~\(R\) with~\(1\) is called a local ring if and only if it has a unique maximal ideal. Suppose that~\(R\) is a ring with a proper ideal~\(M\). Show that~\(R\) is local, with maximal ideal~\(M\), if and only if every element of \(R \setminus M\) (set difference) is a unit.
\end{examproblems}
\end{document}
